\chapter{Introduction}
\label{chap:introduction}

\section{Context and Problem}
Infrastructure as Code makes deployment logic reviewable and repeatable, but
Ansible repositories mix YAML structure, Jinja expressions, inventory values,
role conventions, dynamic includes, runtime facts, and external collections.
A location-only linter therefore cannot always distinguish a defect from a
value defined on another valid execution path. Conversely, a failed provider
must not be interpreted as evidence that a repository is clean.

Existing tools cover important parts of this problem. ansible-lint checks
Ansible syntax and established best practices; GitLeaks searches text and
history for secret patterns. Neither is replaced by this work. The missing
piece is a repository-level semantic and evidence layer that preserves each
provider's result, reasons about Ansible-specific scope, and emits a stable
contract suitable for comparisons and research analysis.

\section{Objectives and Research Questions}
The internship objective was to create that layer from first principles and
evaluate it without overstating what the collected data proves. Four research
questions guide the report:
\begin{description}
  \item[RQ1 -- Coverage.] Which review-checklist concerns and Ansible quality
    mechanisms can be automated, approximated, or must remain manual?
  \item[RQ2 -- Reliability.] Can frozen repositories be scanned while
    preserving partial provider runs and parser failures as diagnostics?
  \item[RQ3 -- Evolution.] Which semantic changes are visible in controlled
    paired runs, including the v0.0.7 role-contract correction?
  \item[RQ4 -- Variable semantics.] Why does \code{unguarded-variable} fire,
    and which cases are unresolved, uncertain, ordered conflicts, or protected
    by definitions and guards?
\end{description}

\section{Contributions}
This work contributes:
\begin{enumerate}
  \item a traceable mapping of 38 review criteria: 15 fully automated, 18
    partially represented, and 5 outside static-analysis scope;
  \item a Python pipeline with tolerant YAML/Jinja parsing, artifact
    classification, internal rules, and repository symbol/execution graphs;
  \item isolated integration of ansible-lint and GitLeaks, including explicit
    completed, partial, unavailable, and failed provider outcomes;
  \item canonical occurrence and provider-evidence models, a packaged
    output-v3 JSON Schema, and full/local/commit/PR workflows;
  \item a reproducible 149-repository historical evaluation, a current
    207-repository diagnostic run, a paired role-contract comparison, and a
    complete numerical-threshold benchmark with a frozen review protocol.
\end{enumerate}

\begin{figure}[htbp]
\centering
\begin{tikzpicture}[
  box/.style={draw=ansblue,rounded corners=3pt,fill=anspaper,
    text width=3.55cm,minimum height=1.35cm,align=center,font=\small},
  arrow/.style={-{Latex[length=2.2mm]},thick,ansblue},
  node distance=1.0cm and .45cm]
  \node[box] (design) {\textbf{Design}\\38 checklist criteria\\76 analysis concepts};
  \node[box,right=of design] (system) {\textbf{Implementation}\\parsing, semantic analysis\\and provider evidence};
  \node[box,right=of system] (study) {\textbf{Evaluation}\\frozen repositories\\paired runs and review};
  \draw[arrow] (design)--(system);
  \draw[arrow] (system)--(study);
  \node[font=\footnotesize,text=ansgray,align=center,
    below=5mm of system] {Observed occurrences are detector evidence; independent human labels are a separate validation step.};
\end{tikzpicture}
\caption{The report's design, implementation, and evaluation path. The numbers describe different units and must not be combined.}
\label{fig:report_roadmap}
\end{figure}

\section{Scope and Claim Boundary}
The reported numbers describe the frozen study population and recorded commit
SHAs. Canonical occurrences are static-analysis evidence, not automatically
human-confirmed defects. \code{CONFIRMED} is a detector status, not a manual
true-positive label. Because a complete, independently adjudicated label set
and a seeded-defect corpus are absent, this report does not claim general
precision or recall. The single-agent diagnostic ratios in
Chapter~\ref{chap:v007} are explicitly limited to their sampled units.

\section{Evidence Hierarchy and Version Policy}
The report follows one evidence hierarchy. Implementation facts are supported
by source code, schemas and automated tests. Execution facts are supported by
logs, health records and frozen outputs. Distribution claims are supported by
the frozen v0.0.6 and v0.0.7 aggregations with their distinct populations.
Mechanism explanations additionally require source sampling. Human accuracy
claims require an independent labelled
benchmark and therefore remain future work.

The current empirical focus is v0.0.7. The v0.0.6 run remains a fully
documented historical study, with v0.0.5 retained only for its paired
same-population comparison. Different repository and metric populations are
labelled and are never combined into one accuracy estimate.

\section{Report Structure}
Chapter~\ref{chap:literature} reviews the research and tool landscape.
Chapter~\ref{chap:methodology} defines the construction and evaluation method.
Chapter~\ref{chap:implementation} documents the implementation and verification
strategy. Chapter~\ref{chap:results} preserves the v0.0.6 study and controlled
v0.0.5 comparator. Chapter~\ref{chap:v007} presents current execution,
variable-resolution review, and threshold-benchmark evidence.
Chapter~\ref{chap:discussion} interprets that evidence and its validity limits.
Chapter~\ref{chap:conclusion} answers the research questions and defines the
next validation work. Appendices~\ref{chap:rules}--\ref{chap:users_outputs}
retain the complete rule catalogue, operational and reproducibility details,
and the field-level output guide.
